% Generated draft: requires mathematical and source review.
\documentclass[11pt]{article}
\usepackage[margin=25mm]{geometry}
\usepackage{fontspec}
\setmainfont{Latin Modern Roman}
\usepackage{amsmath,amssymb,mathtools,unicode-math}
\setmathfont{Latin Modern Math}
\usepackage{enumitem,needspace,xurl,hyperref}
\setlength{\parindent}{0pt}
\setlength{\parskip}{5pt}
\setlength{\emergencystretch}{3em}
\newcommand{\N}{\mathbb{N}}
\newcommand{\Z}{\mathbb{Z}}
\newcommand{\Q}{\mathbb{Q}}
\newcommand{\R}{\mathbb{R}}
\newcommand{\C}{\mathbb{C}}
\newcommand{\F}{\mathbb{F}}
\newcommand{\A}{\mathbb{A}}
\newcommand{\PP}{\mathbb{P}}
\newcommand{\E}{\mathbb{E}}
\newcommand{\eps}{\varepsilon}
\newcommand{\dd}{\,\mathrm{d}}
\begin{document}
\section*{UnsolvedMath ID: 30001121}
\textbf{Source problem code:} catalogue entry
\subsection{Definitions and mathematical statement}
Fix arbitrary parameters \(\gamma>1\) and \(\rho_0>0\). Let \(\Lambda\subset\R^2\) be the domain described in the supplied record as the upper half-plane outside a wedge. The wedge angle, orientation, vertex, and precise domain are not specified. Write \(\xi=(\xi_1,\xi_2)\) for the self-similar spatial variables, and let \(D\) and \(\operatorname{div}\) denote differentiation with respect to \(\xi\). Let \(\Gamma_{\mathrm w}\) denote the wedge portion of \(\partial\Lambda\), and let \(\nu\) be a unit normal there wherever defined.

For a real-valued potential \(\phi\), set
\[
H_\phi=\rho_0-(\gamma-1)\left(\phi+\frac12|D\phi|^2\right),
\qquad
\rho_\phi=\rho(|D\phi|^2,\phi)=H_\phi^{1/(\gamma-1)}.
\]
The physical, real, nonnegative-density interpretation requires \(H_\phi\geq0\) wherever this expression is used. Whether the intended problem requires strictly positive density or permits vacuum is not specified.

\textbf{Requested existence problem.} Seek a self-similar weak solution \(\phi\) on \(\Lambda\) satisfying
\[
\operatorname{div}(\rho_\phi D\phi)+2\rho_\phi=0
\quad\text{in }\Lambda,
\qquad
D\phi\cdot\nu=0
\quad\text{on }\Gamma_{\mathrm w},
\]
and approaching at infinity the two incident constant-state potentials separated by the incident shock. Those potentials, their parameters, the incident shock, and the meaning of this far-field convergence are not specified in the supplied record.

At the level of the interior equation alone, its distributional interpretation is
\[
\int_\Lambda\left(\rho_\phi D\phi\cdot D\eta-2\rho_\phi\eta\right)\,\dd\xi=0
\qquad\text{for every }\eta\in C_c^\infty(\Lambda),
\]
provided \(\rho_\phi\) and \(\rho_\phi D\phi\) are locally integrable. This identity does not specify the intended solution regularity, boundary traces, or shock admissibility. The record also does not state whether a boundary condition is required on any remaining portion of the boundary of the upper half-plane. Consequently, the supplied formulation describes an existence problem but does not determine a fully specified boundary-value problem or its admissible parameter range.
\subsection{Short English statement}
For given \(\gamma>1\) and \(\rho_0>0\), seek a self-similar weak potential-flow solution outside a wedge in the upper half-plane. It must satisfy the stated density law and equation, obey slip on the wedge, and approach two incident constant states separated by a shock at infinity. The wedge geometry, incident data, admissibility conditions, and precise far-field requirement remain unspecified and require mathematical review.
\subsection{Sources}
\begin{itemize}
\item Local UnsolvedMath registry: \texttt{data/unsolvedmath/problems.json}, with \texttt{catalog/problems.json} as the additional-record registry; supplied permanent ID \(30001121\), exact source problem code \texttt{catalogue entry}, title \emph{Self-Similar Shock Diffraction by Wedges}. The supplied record belongs to the Oberwolfach Reports Open Problems collection and gives proposed year \(2008\). No individual proposer is named.
\item Supplied report citation: \emph{Hyperbolic Conservation Laws} (2009), Oberwolfach Reports, pp.\ 3139--3200, DOI \texttt{10.4171/owr/2008/56}. Source URL and literature-source link labeled ``original OWR report'': \url{https://doi.org/10.4171/owr/2008/56}. The report is described as documenting workshop discussions and open problems. The supplied year label \(2009\), proposed year \(2008\), and DOI containing \(2008\) are retained without resolving their bibliographic relationship.
\item Supplied literature-source link labeled ``Chen--Feldman shock diffraction'': \url{https://doi.org/10.4007/annals.2010.171.1067}. No further bibliographic details are supplied.
\item Literature assessment dated exactly \texttt{2026-08-21}: this is supplied dated curation, not independently verified. It labels the problem partially solved and reports global self-similar shock-diffraction solutions for potential flow for large-angle wedges, including existence and stability, while describing the full parameter range and general formulation as not completely settled. The background expressly requests independent verification of time-sensitive status and cited claims.
\item The supplied statement status is \texttt{corrected\_verified}; its verification note says that the mathematical content was checked against the official Oberwolfach report and that notation, context, and question form were normalized without changing the source-level claim. This is a supplied verification assertion, not an independent verification performed for this draft.
\end{itemize}
\end{document}
